% GENERATED FILE. Edit canonical lessons, then run npm run generate:books.
% canonical-source-hash: fnv1a64:dfccc7f4c3fb9957
% canonical-lessons: ML-C462-risarvesan, ML-C462-berttu, ML-C462-kamparttmenru, ML-C462-bukku-ceyyuka, ML-C462-tikkarretukkuka

\chapter{Booking a Berth}
\label{ch:ml-booking-a-berth}
\hlchaptermodality{462}

\begin{chapteropening}
\textbf{By the end of this chapter:} \emph{I can say a reservation, a berth, a compartment, to make a booking and to buy a ticket.}

Complete the last lesson of chapter 462: I can say a reservation, a berth, a compartment, to make a booking and to buy a ticket.
\end{chapteropening}

\section[\texorpdfstring{\ml{റിസർവേഷൻ} (ṟisarvēṣan)}{റിസർവേഷൻ (ṟisarvēṣan)}]{\ml{റിസർവേഷൻ} (ṟisarvēṣan) — a reservation}

\label{lesson:ML-C462-risarvesan}



\begin{quote}
Before the new one: say the Malayalam for a queue, then the Malayalam for to travel.
\end{quote}

\subsection*{You'll want to know: \ml{റിസർവേഷൻ}}
\textbf{\ml{റിസർവേഷൻ}} — \emph{ṟisarvēṣan} — \textquotedblleft{}a reservation\textquotedblright{}.

\begin{quote}
\ml{റിസർവേഷൻ} \ml{ഇല്ലെങ്കിൽ} \ml{സീറ്റ്} \ml{കിട്ടില്ല}.
\end{quote}

\begin{quote}
\emph{ṟisarvēṣan illeṅkil sīṟṟŭ kiṭṭilla.}
\end{quote}

\begin{quote}
If you have no reservation, you won't get a seat.
\end{quote}

A reserved seat on a train is a \ml{റിസർവേഷൻ} \ml{സീറ്റ്}.

\begin{cousinweb}
\textbf{\ml{റിസർവേഷൻ}} (\emph{ṟisarvēṣan}) was borrowed from English \emph{reservation}, a seat or room kept for someone in advance.
\end{cousinweb}

\begin{grammarlens}[title={What you've built}]
One more word for getting around.
\end{grammarlens}

\subsection*{Your turn}
\begin{itemize}
\raggedright
  \item \emph{Say it:} \emph{ṟisarvēṣan}
  \item \emph{Say it:} \emph{ṟisarvēṣan}, once more
  \item \emph{Say it:} \emph{ṟisarvēṣan illeṅkil sīṟṟŭ kiṭṭilla.} — If you have no reservation, you won't get a seat.
  \item \emph{Recall:} say the Malayalam for a queue, then the Malayalam for to travel, then say \emph{ṟisarvēṣan} again
\end{itemize}

\begin{tcolorbox}[breakable,colback=teal!4,colframe=teal!35!black,title={Before you move on}]
What does \ml{റിസർവേഷൻ} mean? (\textquotedblleft{}A reservation\textquotedblright{}.) Say it once more.
\end{tcolorbox}
\section[\texorpdfstring{\ml{ബെർത്ത്} (berttŭ)}{ബെർത്ത് (berttŭ)}]{\ml{ബെർത്ത്} (berttŭ) — a berth}

\label{lesson:ML-C462-berttu}



\begin{quote}
Before the new one: say the Malayalam for to travel, then the Malayalam for a reservation.
\end{quote}

\subsection*{You'll want to know: \ml{ബെർത്ത്}}
\textbf{\ml{ബെർത്ത്}} — \emph{berttŭ} — \textquotedblleft{}a berth\textquotedblright{}.

\begin{quote}
\ml{എനിക്ക്} \ml{മുകളിലെ} \ml{ബെർത്ത്} \ml{ആണ്} \ml{കിട്ടിയത്}.
\end{quote}

\begin{quote}
\emph{enikkŭ mukaḷile berttŭ āṇŭ kiṭṭiyatŭ.}
\end{quote}

\begin{quote}
It was the upper berth that I got.
\end{quote}

The lower berth is \ml{താഴത്തെ} \ml{ബെർത്ത്}; the one in the middle is \ml{നടുവിലെ} \ml{ബെർത്ത്}.

\begin{cousinweb}
\textbf{\ml{ബെർത്ത്}} (\emph{berttŭ}) was borrowed from English \emph{berth}, a bed on a train or a ship.
\end{cousinweb}

\begin{grammarlens}[title={What you've built}]
One more word for getting around.
\end{grammarlens}

\subsection*{Your turn}
\begin{itemize}
\raggedright
  \item \emph{Say it:} \emph{berttŭ}
  \item \emph{Say it:} \emph{berttŭ}, once more
  \item \emph{Say it:} \emph{enikkŭ mukaḷile berttŭ āṇŭ kiṭṭiyatŭ.} — It was the upper berth that I got.
  \item \emph{Recall:} say the Malayalam for to travel, then the Malayalam for a reservation, then say \emph{berttŭ} again
\end{itemize}

\begin{tcolorbox}[breakable,colback=teal!4,colframe=teal!35!black,title={Before you move on}]
What does \ml{ബെർത്ത്} mean? (\textquotedblleft{}A berth\textquotedblright{}.) Say it once more.
\end{tcolorbox}
\section[\texorpdfstring{\ml{കമ്പാർട്ട്മെന്റ്} (kampārṭṭmenṟŭ)}{കമ്പാർട്ട്മെന്റ് (kampārṭṭmenṟŭ)}]{\ml{കമ്പാർട്ട്മെന്റ്} (kampārṭṭmenṟŭ) — a compartment}

\label{lesson:ML-C462-kamparttmenru}



\begin{quote}
Before the new one: say the Malayalam for a reservation, then the Malayalam for a berth.
\end{quote}

\subsection*{You'll want to know: \ml{കമ്പാർട്ട്മെന്റ്}}
\textbf{\ml{കമ്പാർട്ട്മെന്റ്}} — \emph{kampārṭṭmenṟŭ} — \textquotedblleft{}a compartment\textquotedblright{}.

\begin{quote}
\ml{ഞങ്ങളുടെ} \ml{കമ്പാർട്ട്മെന്റിൽ} \ml{നല്ല} \ml{തിരക്കായിരുന്നു}.
\end{quote}

\begin{quote}
\emph{ñaṅṅaḷuṭe kampārṭṭmenṟil nalla tirakkāyirunnu.}
\end{quote}

\begin{quote}
Our compartment was very crowded.
\end{quote}

A train is a row of compartments; a sleeper compartment is a \ml{സ്ലീപ്പർ} \ml{കമ്പാർട്ട്മെന്റ്}.

\begin{cousinweb}
\textbf{\ml{കമ്പാർട്ട്മെന്റ്}} (\emph{kampārṭṭmenṟŭ}) was borrowed from English \emph{compartment}, one of the sections of a railway carriage.
\end{cousinweb}

\begin{grammarlens}[title={What you've built}]
One more word for getting around.
\end{grammarlens}

\subsection*{Your turn}
\begin{itemize}
\raggedright
  \item \emph{Say it:} \emph{kampārṭṭmenṟŭ}
  \item \emph{Say it:} \emph{kampārṭṭmenṟŭ}, once more
  \item \emph{Say it:} \emph{ñaṅṅaḷuṭe kampārṭṭmenṟil nalla tirakkāyirunnu.} — Our compartment was very crowded.
  \item \emph{Recall:} say the Malayalam for a reservation, then the Malayalam for a berth, then say \emph{kampārṭṭmenṟŭ} again
\end{itemize}

\begin{tcolorbox}[breakable,colback=teal!4,colframe=teal!35!black,title={Before you move on}]
What does \ml{കമ്പാർട്ട്മെന്റ്} mean? (\textquotedblleft{}A compartment\textquotedblright{}.) Say it once more.
\end{tcolorbox}
\section[\texorpdfstring{\ml{ബുക്ക്} \ml{ചെയ്യുക} (bukkŭ ceyyuka)}{ബുക്ക് ചെയ്യുക (bukkŭ ceyyuka)}]{\ml{ബുക്ക്} \ml{ചെയ്യുക} (bukkŭ ceyyuka) — to make a booking; to book}

\label{lesson:ML-C462-bukku-ceyyuka}



\begin{quote}
Before the new one: say the Malayalam for a berth, then the Malayalam for a compartment.
\end{quote}

\subsection*{You'll want to know: \ml{ബുക്ക്} \ml{ചെയ്യുക}}
\textbf{\ml{ബുക്ക്} \ml{ചെയ്യുക}} — \emph{bukkŭ ceyyuka} — \textquotedblleft{}to make a booking; to book\textquotedblright{}.

\begin{quote}
\ml{അച്ഛൻ} \ml{ഓൺലൈനായി} \ml{രണ്ട്} \ml{ടിക്കറ്റ്} \ml{ബുക്ക്} \ml{ചെയ്തു}.
\end{quote}

\begin{quote}
\emph{acchan ōṇlaināyi raṇṭŭ ṭikkaṟṟŭ bukkŭ ceytu.}
\end{quote}

\begin{quote}
Father booked two tickets online.
\end{quote}

Past: \ml{ബുക്ക്} \ml{ചെയ്തു}, \emph{bukkŭ ceytu}.

English book plus \ml{ചെയ്യുക}, to do. Malayalam makes many new verbs this way, as in \ml{പാർക്ക്} \ml{ചെയ്യുക}, to park.

\begin{cousinweb}
\textbf{\ml{ബുക്ക്} \ml{ചെയ്യുക}} (\emph{bukkŭ ceyyuka}) is a Malayalam verb phrase built from \ml{ബുക്ക്}, which was borrowed from English \emph{book}, \textquotedblleft{}to reserve\textquotedblright{}, and the native verb \ml{ചെയ്യുക}, to do.
\end{cousinweb}

\begin{grammarlens}[title={What you've built}]
One more word for getting around.
\end{grammarlens}

\subsection*{Your turn}
\begin{itemize}
\raggedright
  \item \emph{Say it:} \emph{bukkŭ ceyyuka}
  \item \emph{Say it:} \emph{bukkŭ ceyyuka}, once more
  \item \emph{Say it:} \emph{acchan ōṇlaināyi raṇṭŭ ṭikkaṟṟŭ bukkŭ ceytu.} — Father booked two tickets online.
  \item \emph{Recall:} say the Malayalam for a berth, then the Malayalam for a compartment, then say \emph{bukkŭ ceyyuka} again
\end{itemize}

\begin{tcolorbox}[breakable,colback=teal!4,colframe=teal!35!black,title={Before you move on}]
What does \ml{ബുക്ക്} \ml{ചെയ്യുക} mean? (\textquotedblleft{}To make a booking; to book\textquotedblright{}.) Say it once more.
\end{tcolorbox}
\section[\texorpdfstring{\ml{ടിക്കറ്റെടുക്കുക} (ṭikkaṟṟeṭukkuka)}{ടിക്കറ്റെടുക്കുക (ṭikkaṟṟeṭukkuka)}]{\ml{ടിക്കറ്റെടുക്കുക} (ṭikkaṟṟeṭukkuka) — to buy a ticket}

\label{lesson:ML-C462-tikkarretukkuka}



\begin{quote}
Before the new one: say the Malayalam for a compartment, then the Malayalam for to make a booking.
\end{quote}

\subsection*{You'll want to know: \ml{ടിക്കറ്റെടുക്കുക}}
\textbf{\ml{ടിക്കറ്റെടുക്കുക}} — \emph{ṭikkaṟṟeṭukkuka} — \textquotedblleft{}to buy a ticket\textquotedblright{}.

\begin{quote}
\ml{ഞാൻ} \ml{സ്റ്റേഷനിൽ} \ml{പോയി} \ml{ടിക്കറ്റെടുത്തു}.
\end{quote}

\begin{quote}
\emph{ñān sṟṟēṣanil pōyi ṭikkaṟṟeṭuttu.}
\end{quote}

\begin{quote}
I went to the station and bought a ticket.
\end{quote}

Past: \ml{ടിക്കറ്റെടുത്തു}, \emph{ṭikkaṟṟeṭuttu}.

\ml{ടിക്കറ്റ്} plus \ml{എടുക്കുക}, to take: in Malayalam you take a ticket. The two words often run together in writing.

\begin{cousinweb}
\textbf{\ml{ടിക്കറ്റെടുക്കുക}} (\emph{ṭikkaṟṟeṭukkuka}) is a Malayalam compound verb built from \ml{ടിക്കറ്റ്}, which was borrowed from English \emph{ticket}, and the native verb \ml{എടുക്കുക}, to take (DEDR 851).
\end{cousinweb}

\begin{grammarlens}[title={What you've built}]
One more word for getting around.
\end{grammarlens}

\subsection*{Your turn}
\begin{itemize}
\raggedright
  \item \emph{Say it:} \emph{ṭikkaṟṟeṭukkuka}
  \item \emph{Say it:} \emph{ṭikkaṟṟeṭukkuka}, once more
  \item \emph{Say it:} \emph{ñān sṟṟēṣanil pōyi ṭikkaṟṟeṭuttu.} — I went to the station and bought a ticket.
  \item \emph{Recall:} say the Malayalam for a compartment, then the Malayalam for to make a booking, then say \emph{ṭikkaṟṟeṭukkuka} again
\end{itemize}

\begin{tcolorbox}[breakable,colback=teal!4,colframe=teal!35!black,title={Before you move on}]
What does \ml{ടിക്കറ്റെടുക്കുക} mean? (\textquotedblleft{}To buy a ticket\textquotedblright{}.) Say it once more.
\end{tcolorbox}
